\documentclass[11pt]{article}
\usepackage[margin=1in]{geometry}
\usepackage{amsmath,amssymb,amsthm}
\usepackage{mathtools}
\usepackage{enumitem}
\usepackage{hyperref}
\usepackage{cleveref}
\usepackage{booktabs}
\usepackage{tabularx}
\usepackage{microtype}
\usepackage{cite}
\usepackage{xurl}
\usepackage{path}
\hypersetup{colorlinks=true,linkcolor=blue,citecolor=blue,urlcolor=blue}
\newtheorem{definition}{Definition}
\newtheorem{lemma}{Lemma}
\newtheorem{remark}{Remark}
\title{Public Request Notice Normal Forms for Source-Only Releases\\\large Canonical Sentence Families for Within-Release Status and Cross-Release Carryforward Notices}
\author{Anonymous}
\date{Unified archive note v0.07 (2026-03-21; archive v473)}
\begin{document}
\maketitle

\begin{abstract}
The archive can now say what a source-only paper promises on request, what omitted-support families that promise covers, when those families are complete, what the paper-level public token is, whether that token survives or reopens under a successor release, and what exact outward-facing sentence readers saw. One smaller publication question still recurs in terse papers and release notes: which sentence family is canonical for that visible notice, so the wrapper sentence is not reinvented each time? This note gives that owner one home. It defines compact \emph{public request notice normal forms}: maintained template-and-slot objects that factor one visible public-request notice into a canonical sentence family plus explicit basis-token bindings, without re-owning the notice instance or the narrower status and carryforward semantics it compresses.
\end{abstract}

\paragraph{Non-redundancy note.}
Synthesis~59 owns the within-release paper-level public status token, Synthesis~60 owns the cross-release carryforward verdict, Synthesis~61 owns the emitted outward-facing notice sentence and reader-pointer order, Synthesis~63 owns the token-keyed choice of which canonical notice family applies, Synthesis~64 owns the smallest-sufficient maintained follow-up handoff once that family choice is fixed, Synthesis~73 owns the request-side terminal stop later terse papers should usually cite, Synthesis~74--75 own the paired terminal / cutover discipline, and Synthesis~17 owns the concrete worked instantiation.\cite{series:publicrequeststatusenvelopes,series:publicrequestcarryforward,series:publicrequestnotices,series:publicrequestnoticeselections,series:publicrequestresponsemenus,series:publicrequestresponsepacketrefreshresponsepacketclosureverdicts,series:pairedterminalcitationdiscipline,series:pairedterminalcutoverrule,series:workedexample}
This note owns only the canonical sentence family and slot realization rule for those already-owned notice objects. Later terse papers should cite it directly only when the exact reusable notice family or slot map is itself live: stop at Synthesis~59--60 when the live issue is the imported status or carryforward token, stop at Synthesis~61 when the live issue is the exact emitted sentence or pointer order, widen to Synthesis~63 when the live issue is the token-keyed family choice, widen through Synthesis~64--72 when the live issue has already moved into the exact follow-up handoff or successor refresh branch, widen to Synthesis~73 when the live issue is whether that visible request-side branch already closes, widen to Synthesis~74--75 when the remaining issue is paired-terminal reuse or fail-closed cutover, and stop at Synthesis~17 when the concrete maintained worked route is the live issue instead.

\paragraph{Scope.}
This is not a new public request contract, not a new completion criterion, not a new paper-level token, not a new carryforward verdict, and not a replacement for the emitted notice.
It is a sentence-family note for the exact short notice form that later papers may reuse once the narrower owners already exist.

\paragraph{Exports.}
This note exports (i) public request notice normal forms, (ii) a basis-compatibility rule, (iii) a slot-realization rule, (iv) an emitted-notice inheritance rule separating one canonical sentence family from the concrete notice instance that realizes it, and (v) a forced family-realization matrix for later terse papers that need the exact reusable token-to-slot-to-sentence surface. If later notes instead need the token-keyed choice of which such family applies, they should import Synthesis~63 rather than widen this note. If they instead need the smallest-sufficient maintained follow-up owner that should be handed over after one such family choice is fixed, they should import Synthesis~64 rather than widen this note. If they need only the stable request-side stop or the stop-versus-widen judgment for the adjacent visible source-only branch, they should cite Synthesis~73--75 rather than widen this note.

\section{Why notice normal forms deserve their own note}
A very short source-only paper often wants to reuse one sentence family across base and successor releases: one template for within-release status, another for carryforward.
Without a maintained normal-form owner, neighboring notes start to cite one concrete sentence as if it were the canonical form, or they silently drift in wording while claiming to say the same thing.
That is a different question from who owns the notice instance itself.
The notice instance owns the emitted sentence and pointer order for one release cut.
The normal form owns the canonical family that the emitted sentence is allowed to realize.

\begin{definition}[Public request notice normal form]
Fix one public request notice $N$ whose basis is either a within-release public request status envelope or a successor-release public request carryforward envelope.
A \emph{public request notice normal form} is a tuple
\[
F=(\tau,\Theta,\Lambda,\rho,\sigma)
\]
consisting of one notice kind $\tau$, one allowed token-compatibility set $\Theta$, one ordered literal-and-slot template $\Lambda$, one realization map $\rho$ sending each slot to an already-owned field in the imported notice or its basis object, and one rendered sentence $\sigma$ produced by substituting those slot values into $\Lambda$.
The normal form does not own the notice instance.
It owns only the canonical family and the explicit slot map that makes the emitted sentence auditable.
\end{definition}

\begin{definition}[Basis-compatibility rule]
A public request notice normal form satisfies the \emph{basis-compatibility rule} if its notice kind and allowed token set agree with the imported emitted notice and its basis object.
For a within-release form, the allowed status token must match the imported public request status envelope.
For a successor-carryforward form, the allowed status token and carryforward token must match the imported carryforward envelope.
So the normal form may not authorize a sentence family for a token combination the basis object does not actually certify.
\end{definition}

\begin{definition}[Slot-realization rule]
A public request notice normal form satisfies the \emph{slot-realization rule} if every non-literal semantic phrase in its rendered sentence is carried by an explicit slot whose source field is named in the realization map, and if substituting those slot values into the template produces the emitted notice sentence exactly.
So the normal form may compress wording, but it may not hide semantics in undeclared glue text.
\end{definition}

\begin{definition}[Emitted-notice inheritance rule]
A public request notice normal form satisfies the \emph{emitted-notice inheritance rule} if every emitted notice key it names already exists in the companion public-request-notice ledger and if the normal form changes neither the reader-pointer order nor the narrower owner artifacts that emitted notice already imports.
So the normal form may factor the sentence family, but it may not replace the concrete notice instance.
\end{definition}

\begin{lemma}[Canonical notice families reduce wording drift without replacing notice ownership]
If a public request notice normal form satisfies the basis-compatibility, slot-realization, and emitted-notice inheritance rules, then a later terse paper may cite the normal form when it needs the canonical short sentence family for a source-only notice, while still citing the emitted notice when it needs to say which concrete sentence and reader pointers were actually shipped.
\end{lemma}

\begin{proof}
The basis-compatibility rule prevents the normal form from licensing a status or carryforward combination that the basis object does not certify.
The slot-realization rule ensures that the rendered sentence is reconstructed exactly from declared literals and owned token fields rather than from hidden prose drift.
The emitted-notice inheritance rule prevents the normal form from silently taking over the concrete notice instance or its reader-pointer order.
Therefore the normal form contributes one reusable sentence family, not a new status owner and not a replacement for the emitted notice.\qedhere
\end{proof}

\begin{remark}[Minimal citation rule]
Cite the public request notice normal form when the claim is about \emph{which canonical sentence family or template was used}.
Cite the public request notice when the claim is about \emph{which exact outward-facing sentence and pointer order were actually emitted}.
Cite the status or carryforward envelope when the claim is about \emph{which token or cross-release verdict is true}.
\end{remark}

\begin{remark}[Normal-form stop discipline]
When later terse papers need only the stable request-side stop or the stop-versus-widen judgment for the visible source-only branch, they should usually cite Synthesis~73--75 rather than expand this family layer locally.
Widen through this note only when the exact wording family or slot map is itself disputed. In that case, the first sufficient audit stop is the selected normal form itself, and widening should proceed only to the emitted public-request notice that realizes it and then to the imported status or carryforward basis object named in that notice.
So the normal form owns the reusable sentence family, while the emitted sentence instance and the narrower basis semantics remain with their already-maintained owners.
\end{remark}

\paragraph{A forced public notice-family matrix.}
\begin{table}[t]
\centering
\small
\begin{tabularx}{0.97\linewidth}{@{}p{0.31\linewidth}X@{}}
\toprule
\textbf{Forced row} & \textbf{Forced value}\\
\midrule
Within-release force row & \texttt{public\_request\_notice\_normal\_form\_key = worked\_example\_receipt\_interlock\_within\_release\_complete\_public\_request\_notice\_normal\_form} with \texttt{notice\_kind = within\_release}, \texttt{basis\_key = worked\_example\_receipt\_interlock\_public\_request\_status\_envelope}, and the lone compatible slot bundle \texttt{public\_request\_status\_token = complete} \\
Successor force row & \texttt{public\_request\_notice\_normal\_form\_key = worked\_example\_receipt\_interlock\_successor\_preserved\_complete\_public\_request\_notice\_normal\_form} with \texttt{notice\_kind = successor\_carryforward}, \texttt{basis\_key = worked\_example\_receipt\_interlock\_public\_request\_carryforward\_envelope}, and the only worked compatible token pair \texttt{public\_request\_status\_token = complete} plus \texttt{carryforward\_token = preserved\_complete} \\
Compatibility gate & Any later terse paper may quote one matrix row only while the imported notice kind and compatible token bundle still match that row exactly; a drift to \texttt{advanced}, \texttt{reopened}, or any different basis object forces a different family or a fail-closed reopen \\
Slot-source floor & The within-release row binds \texttt{public\_request\_status\_token} from \nolinkurl{example_public_request_status_envelope.json::public_request_status_token}; the successor row binds \texttt{public\_request\_status\_token} from \nolinkurl{example_public_request_carryforward_envelope.json::successor_public_request_status_token} and binds \texttt{carryforward\_phrase} only through the explicit render map \texttt{preserved\_complete -> unchanged} \\
Rendered-sentence floor & Within release the forced rendered sentence is ``For this source-only release, the worked note's paper-level public request promise is complete and supporting artifacts remain available on request.'' Across the successor pair the forced rendered sentence is ``Across the canned successor release, that complete paper-level public request status carries forward unchanged.'' \\
Inherited-wrapper floor & The within-release row still inherits \texttt{emitted\_notice\_key = worked\_example\_receipt\_interlock\_public\_request\_status\_notice}; the successor row still inherits \texttt{worked\_example\_receipt\_interlock\_public\_request\_carryforward\_notice}; neither row replaces reader-pointer order or basis ownership \\
Staleness trigger & This matrix goes stale immediately if the compatible token bundle, a named slot source, the render map, or the inherited emitted-notice key changes, because then the quoted sentence family would stop being the same canonical reusable family \\
Escalation pointer & Widen to Synthesis~59--60 for basis-token / carryforward semantics, Synthesis~61 for the exact emitted wrapper, Synthesis~63 for token-to-family choice, and Synthesis~73--75 once the live issue is request-side terminal reuse or cutover \\
\bottomrule
\end{tabularx}
\caption{A forced public notice-family matrix. This is the smallest public answer later terse papers can quote directly when the live issue is the exact reusable source-only sentence family and slot map rather than the narrower basis-token owner, the emitted wrapper instance, or the request-side terminal tail.}
\label{tab:forced-public-notice-family-matrix}
\end{table}

This matrix is intentionally non-decorative: it pins the actual worked within-release and successor-facing normal-form keys, the exact compatible token bundles, the explicit slot sources, the rendered-sentence floor, and the point at which the family citation has gone stale without silently re-owning the emitted notice instance, its reader-pointer order, or the narrower status / carryforward semantics that those wrappers compress.\cite{series:workedexample}

\paragraph{One tiny family-forcing drill.}
In the worked within-release row, the only semantic slot is \texttt{public\_request\_status\_token}, and the basis object fixes it to \texttt{complete}, so the rendered sentence may differ only if that imported token or its named slot source changes. In the worked successor row, the second semantic phrase is not free prose: \texttt{carryforward\_phrase} is bound only through the explicit render map \texttt{preserved\_complete -> unchanged}. If the carryforward envelope instead published \texttt{advanced} or \texttt{reopened}, or if the emitted wrapper stopped inheriting the same notice key, this matrix could no longer be quoted unchanged: the paper would have to cite a different normal form or fail closed until one exists.\cite{series:workedexample}

\section{Worked example pointer}
The maintained worked bundle now ships \nolinkurl{example_public_request_notice_normal_forms.json}.\cite{series:workedexample}
It contains exactly two normal-form entries keyed \nolinkurl{worked_example_receipt_interlock_within_release_complete_public_request_notice_normal_form} and \nolinkurl{worked_example_receipt_interlock_successor_preserved_complete_public_request_notice_normal_form}.
The within-release normal form binds one slot \nolinkurl{public_request_status_token} back to \nolinkurl{example_public_request_status_envelope.json::public_request_status_token} and renders the exact sentence ``For this source-only release, the worked note's paper-level public request promise is complete and supporting artifacts remain available on request.''
The successor-carryforward normal form binds \nolinkurl{public_request_status_token} to \nolinkurl{example_public_request_carryforward_envelope.json::successor_public_request_status_token} and binds \nolinkurl{carryforward_phrase} through the explicit render map \nolinkurl{preserved_complete -> unchanged}, yielding the exact sentence ``Across the canned successor release, that complete paper-level public request status carries forward unchanged.''
That gives later terse papers one citeable owner for the canonical sentence family without confusing that family with the concrete emitted notice. In the maintained worked bundle, those two exact normal-form keys are imported directly onto the shipped audience/question routes beside the corresponding emitted notice keys, so the canonical wrapper family is recoverable from one maintained route object without re-owning reader-pointer order. Later terse papers should usually combine that route-visible family layer with Synthesis~73--75 rather than restitch the visible source-only branch locally, widen through Synthesis~63 only when the token-to-family choice itself is disputed, and widen through Synthesis~64--72 only when the live issue has already moved into the exact maintained follow-up handoff or successor refresh branch.\cite{series:publicrequestnoticeselections,series:publicrequestresponsemenus,series:publicrequestresponsepackets,series:publicrequestresponsepacketcarryforward,series:publicrequestresponsepacketdeltaledgers,series:publicrequestresponsepacketrefreshnotices,series:publicrequestresponsepacketrefreshnoticenormalforms,series:publicrequestresponsepacketrefreshnoticeselections,series:publicrequestresponsepacketrefreshresponsemenus,series:publicrequestresponsepacketrefreshresponsepackets,series:publicrequestresponsepacketrefreshresponsepacketclosureverdicts,series:pairedterminalcitationdiscipline,series:pairedterminalcutoverrule}

\begin{thebibliography}{99}\small

\bibitem{series:publicrequeststatusenvelopes}
Anonymous.
\newblock Public Request Status Envelopes for Source-Only Releases: Paper-Level Tokens Compressing Family Completion without Re-owning the Promise.
\newblock Series synthesis note (Synthesis~59), 2026. (This archive.)

\bibitem{series:publicrequestcarryforward}
Anonymous.
\newblock Public Request Carryforward for Source-Only Successor Releases: When Paper-Level Request Status Persists, Advances, or Reopens across Release Evolution.
\newblock Series synthesis note (Synthesis~60), 2026. (This archive.)

\bibitem{series:publicrequestnotices}
Anonymous.
\newblock Public Request Notices for Source-Only Releases: Outward-Facing Sentences for Within-Release Status and Cross-Release Carryforward.
\newblock Series synthesis note (Synthesis~61), 2026. (This archive.)

\bibitem{series:publicrequestnoticeselections}
Anonymous.
\newblock Public Request Notice Selections for Source-Only Releases: Choosing Canonical Notice Families from Within-Release Status and Cross-Release Carryforward Tokens.
\newblock Series synthesis note (Synthesis~63), 2026. (This archive.)

\bibitem{series:publicrequestresponsemenus}
Anonymous.
\newblock Public Request Response Menus for Source-Only Releases: Choosing the Smallest Sufficient Maintained Follow-up Object.
\newblock Series synthesis note (Synthesis~64), 2026. (This archive.)

\bibitem{series:publicrequestresponsepackets}
Anonymous.
\newblock Public Request Response Packets for Source-Only Releases: Exact Bounded Follow-up Slices for Shipped Source-Only Notices.
\newblock Series synthesis note (Synthesis~65), 2026. (This archive.)

\bibitem{series:publicrequestresponsepacketcarryforward}
Anonymous.
\newblock Public Request Response Packet Carryforward for Source-Only Successor Releases: Reusing or Refreshing Exact Follow-up Slices.
\newblock Series synthesis note (Synthesis~66), 2026. (This archive.)

\bibitem{series:publicrequestresponsepacketdeltaledgers}
Anonymous.
\newblock Public Request Response Packet Delta Ledgers for Source-Only Successor Releases: Exact Refresh Deltas for Reused Follow-up Slices.
\newblock Series synthesis note (Synthesis~67), 2026. (This archive.)

\bibitem{series:publicrequestresponsepacketrefreshnotices}
Anonymous.
\newblock Public Request Response Packet Refresh Notices for Source-Only Successor Releases: Emitted Visible Reissue Sentences for Refreshed Follow-up Slices.
\newblock Series synthesis note (Synthesis~68), 2026. (This archive.)

\bibitem{series:publicrequestresponsepacketrefreshnoticenormalforms}
Anonymous.
\newblock Public Request Response Packet Refresh Notice Normal Forms for Source-Only Successor Releases: Canonical Visible Wrapper Families for Refreshed Follow-up Slices.
\newblock Series synthesis note (Synthesis~69), 2026. (This archive.)

\bibitem{series:publicrequestresponsepacketrefreshnoticeselections}
Anonymous.
\newblock Public Request Response Packet Refresh Notice Selections for Source-Only Successor Releases: Choosing Visible Wrapper Families for Refreshed Follow-up Slices.
\newblock Series synthesis note (Synthesis~70), 2026. (This archive.)

\bibitem{series:publicrequestresponsepacketrefreshresponsemenus}
Anonymous.
\newblock Public Request Response Packet Refresh Response Menus for Source-Only Successor Releases: Choosing the Smallest Sufficient Successor-Facing Handoff after a Visible Refresh.
\newblock Series synthesis note (Synthesis~71), 2026. (This archive.)

\bibitem{series:publicrequestresponsepacketrefreshresponsepackets}
Anonymous.
\newblock Public Request Response Packet Refresh Response Packets for Source-Only Successor Releases: Exact Successor-Facing Handoff Slices after a Visible Refresh.
\newblock Series synthesis note (Synthesis~72), 2026. (This archive.)


\bibitem{series:publicrequestresponsepacketrefreshresponsepacketclosureverdicts}
Anonymous.
\newblock Public Request Response Packet Refresh Response Packet Closure Verdicts for Source-Only Successor Releases: Capstone Certificates for Exact Successor-Facing Handoff Slices.
\newblock Series synthesis note (Synthesis~73), 2026. (This archive.)

\bibitem{series:pairedterminalcitationdiscipline}
Anonymous.
\newblock Paired Terminal Citation Discipline for Stable Review Spines: One Answer-Side Stop and One Request-Side Capstone for Terse Papers.
\newblock Series synthesis note (Synthesis~74), 2026. (This archive.)

\bibitem{series:pairedterminalcutoverrule}
Anonymous.
\newblock Paired Terminal Cutover Rules for Stable Review Spines: When Later Terse Papers Must Stop, Widen, or Reopen.
\newblock Series synthesis note (Synthesis~75), 2026. (This archive.)

\bibitem{series:workedexample}
Anonymous.
\newblock Worked Example: Receipt Interlock for an Anonymous-DHT Reader-Privacy Claim.
\newblock Series synthesis note (Synthesis~17), 2026. (This archive.)

\end{thebibliography}
\end{document}
